\scanpage{49}
Soient \deleted{\unsure{les}}
\[
 \underset{\substack{\shortparallel\\ L^\perp\otimes L}}{\underline L},\
 \underset{\substack{\shortparallel\\ M^\perp\otimes M}}{\underline M}
 \subset\mathcal M=\underline{\operatorname{End}}(V)
\]
correspondant à $L,M$, \deleted{\unsure{à ceux-ci|ils sont}} \inserted{ils sont} caractérisés par la condition de correspondre à des sections de $X$, i.e.
\[
 \left\{\begin{aligned}
 \underline L\subset\overline{\mathfrak g},\quad
 q\mid\underline L=0&\qquad\text{i.e.}\quad\operatorname{Tr},\det\mid\underline L=0\\
 \underline M\subset\widetilde{\mathfrak g},\quad
 q\mid\underline M=0&\qquad\text{i.e.}\quad\operatorname{Tr},\det\mid\underline M=0
 \end{aligned}\right.\tag*{(8)}
\]
et que de plus on ait, posant
\[
 \left\{\begin{aligned}
 J&=\underline{\mathcal O}_S.u\subset\mathcal M
 \qquad(J\in\Gamma(Q/S))\\
 J\circ L&=0\\
 \underline M\circ J&=0
 \end{aligned}\right.\tag*{(9)}
\]
\noindent\rule{.10\linewidth}{.3pt}\\[-.8em]\rule{.11\linewidth}{.3pt}

Notons que la condition $J\in\Gamma(Q/S)$ implique que
\[
 \mathcal A=\underline{\mathcal O}_S\oplus J
 \underset{\underline{\text{univ}}}{\lhook\joinrel\unsure{-|\longrightarrow}}
 \mathcal M=\underline{\operatorname{End}}(V)\tag*{(10)}
\]
est un ss-fibré \deleted{\unsure{de}} vectoriel de rang 2, \deleted{et les} \deleted{conditions (9)} La donnée de $J$ \uline{équivaut} à la donnée de $\mathcal A$, \uline{et} d'un homomorphisme d'Algèbres
% slanted
\margin{expliciter un peu}
\[
 \mathcal A\xrightarrow{\rho}\underline{\mathcal O}_S\tag*{(11)}
\]
(savoir l'unique hom d'Alg. de noyau $J$).

Ainsi on a un iso
\[
\begin{aligned}
 Q\simeq{}&\text{ens des couples }\mathcal A,\rho,\\
 &\mathcal A\subset\mathcal M\ \text{sous-algèbre }\deleted{\text{\unsure{univ. de rg 2}}}\ \text{\unsure{spéc|spéc.|spéciale}}\\
 &\rho:\mathcal A\longrightarrow\underline{\mathcal O}_S\quad\text{hom d'Alg}
\end{aligned}\tag*{(12)}
\]
% archive sheet number: 48
